% Slide build of Figure 12.1 -- one triangle that evolves, not two side by side.
%
%   1  three switches, three cables: physically redundant, and fine
%   2  one broadcast starts going round, and nothing stops it
%   3  spanning tree blocks one port -- the cable is still there
%
% The cable that gets blocked stays drawn at every level, so Morph recolours
% it rather than removing and re-adding it. That is the point the caption
% makes: spanning tree does not remove redundancy, it holds it in reserve.
\begin{tikzpicture}[font=\scriptsize]
  \tikzset{
    dv/.style={rounded corners=3pt, draw=#1, fill=#1!8, text=#1!45!black,
               line width=1.1pt, minimum width=1.8cm, minimum height=1.2cm,
               align=center, font=\small\bfseries},
    w/.style={line width=1.4pt, neutral!70},
    % text width is not optional: without it the caption is one long line
    % that runs past the pinned bounding box and is cropped.
    capt/.style={font=\bfseries, anchor=north west, text width=9.4cm,
                 align=left}}

  \useasboundingbox (-1.0,-2.6) rectangle (9.6,3.0);

  \node[dv=dom2] (a) at (4.0,2.0) {\faIcon{sitemap}\\SW1};
  \node[dv=dom2] (b) at (1.9,-0.1) {\faIcon{sitemap}\\SW2};
  \node[dv=dom2] (c) at (6.1,-0.1) {\faIcon{sitemap}\\SW3};

  \draw[w] (a) -- (b);
  \draw[w] (a) -- (c);
  % the third cable: drawn at every level, recoloured at level 3 so Morph
  % changes its colour rather than removing and re-adding the shape
  \ifnum\buildlevel<3
    \draw[w] (b) -- (c);
  \else
    \draw[w, accent!70, dash pattern=on 4pt off 3pt] (b) -- (c);
  \fi

  \steponly{1}{\node[capt, text=neutral!45!black] at (-0.6,-1.5)
      {Three switches, three cables. The redundancy is deliberate.};}

  \steponly{2}{%
    % An arc, not a circle: TikZ draws no arrowhead on a closed path, so a
    % \draw[->] ... circle renders a plain ring with no sense of direction.
    \draw[-{Stealth[length=3.2mm]}, accent!85, line width=2.2pt]
        (4.72,0.62) arc (0:308:0.72);
    \node[capt, text=accent!75!black] at (-0.6,-1.5)
        {One broadcast, and an Ethernet frame has no TTL --- it circulates
         for ever.};}

  \steponly{3}{%
    \node[circle, draw=accent!85, fill=white, line width=1.4pt,
          inner sep=2pt, font=\tiny\bfseries, text=accent!75!black]
        at (4.0,-0.1) {BLK};
    \node[font=\small\bfseries, text=dom2!45!black, anchor=west]
        at (5.05,2.0) {ROOT};
    \node[capt, text=dom2!50!black] at (-0.6,-1.5)
        {One port blocks. The cable is still there, ready if another
         link fails.};}
\end{tikzpicture}
